\scanpage{85}

% L1, cancelled opening
\deleted{Soit $C$ $n$-maquette \unsure{rég.}\inserted{régulière}, \unsure{correspond|correspon}}

% L2–7
Soit $\unsure{C|\mathcal C}$ une $n$-carte régulière combinatoire, \deleted{\unsure{repérage}} $R$ l'ens. de \deleted{ses} repères, qui est un torseur à gauche sous $G=\operatorname{Aut}C$, le commutant de $G$ étant formé des opérations du groupe cartographique, engendré par $\sigma_0,\ldots,\sigma_n$. Si $\underline r\in R$, il lui correspond des générateurs $\unsure{\tau_i|\mathrm T_i}=\tau_i^{(\underline r)}$ de $G$, définis par

% L8
\[
\tau_i^{(\underline r)}\underline r=\sigma_i\underline r \tag*{(1)}
\]

% L9–11
Une réalisation géométrique projective de $C$, \unsure{dans|ds} $(n+1)$-fibré projectif $X$ sur $S$, revient à la donnée de a) Opération $\varphi$ de $C$ sur $X$

% L12
\[
\varphi:G\longrightarrow\operatorname{Aut}_S X=\Gamma\mathop{\underline{\operatorname{Aut}}}\nolimits_S(X)\tag*{(2)}
\]

% L13–14
b) drapeau \inserted{max} $r=(\deleted{\unsure{\Phi_0}}X_0,\deleted{\unsure{\Phi_1}}X_1,\ldots,X_{\unsure{n|n+1}})\deleted{X}$ dans $X$ \deleted{\unsure{avec}} satisfaisant

% L15–16
\[
\begin{aligned}
\tau_iX_j&=X_j\quad\text{si } i\ne j\\
\deleted{\tau_iX_i}&\deleted{\ne X_i}
\end{aligned}\tag*{(3)}
\]

% L17
(4) Les $\tau_i$ sont des réflexions de $P$

% L18
[ Ceci permet de \inserted{définir $R\longrightarrow\text{Drap.max}(X)$ et de} factoriser en $\operatorname{Drap}_{\unsure{+1|\ne1|H}}(C)\longrightarrow\operatorname{Drap}_{H}(X)$ ]

% L19
\uline{Prop} Conditions équivalentes (\deleted{\unsure{si on est}} \unsure{$S=\operatorname{Spec}(k)$|$S=\operatorname{Spec}(K)$} \inserted{$k$ corps})

% L20–22
b) Aucune des applications

% L21
\[
\deleted{\unsure{(4)}}\qquad \Phi_i\longrightarrow\operatorname{Grass}_i(X)\qquad 0\leq i\leq n
\]

% L22
n'est constante, \deleted{sur aucune fibre géom}

% L23
\deleted{b$'$) $\forall$ \unsure{$i$}, $0\leq i\leq n$, on a $\tau_i X_i\ne X_i$.}

% L24, left brace
\margin{indépendant du drapeau $r$.}
% slanted

% L24
d) si $\underline f_{i-1}<\underline f_i\ne\underline f'_i<\underline f_{i+1}$ ds $\Phi_*(C)$, alors $f_i\ne f'_i$

% L25–26
a) L'image de $\underline\Phi_0$ engendre lin$^{\mathrm t}$ $X$ (i.e. l'orbite de $X_0$ engendre lin$^{\mathrm t}$ $X$

% L26
\uline{a$'$)} L'intersection des hyperplans $\in\Phi_n$ est vide

% L27
\deleted{\unsure{d)}} c) $\forall\ 0\leq i\leq n$, $\tau_iX_i\ne X_i$

% archive sheet number: 84
